% 08-market-opportunity.tex - Three-segment market sizing (three units, never summed)
%
% STRUCTURAL RULE: S1 is counted in INDIVIDUALS, S2 in INSTITUTIONS and S3 in
% COMPANIES. Those counts are never added together. Only annual platform
% revenue in dollars is summed across segments.
%
% Every figure below is either (a) expanded from a macro in values/variables.tex
% or (b) carries a citation. Assumptions that research could NOT source are
% marked with a TODO(SOURCE) comment immediately above the sentence that uses
% them, and are also flagged in the visible text.

\section{Market Opportunity}

The previous version of this section summed four heterogeneous populations---%
Bitcoin holders, password-manager users, physical-vault owners and private-security
customers---into a single \numint{9000000}-person ``subscriber'' market. That sum was
dimensionally meaningless: it added people who would buy four different products, at
four different prices, through four different channels. It has been replaced by three
segments, each with its own unit of count, its own ARPU and its own penetration
assumption. \textbf{The three headcounts are never added together.} The only quantity
that is legitimately comparable across them is annual platform revenue in dollars.

\subsection{Segment Structure}

\begin{table}[H]
\centering
\begin{tabularx}{\linewidth}{l l Y Y Y}
\toprule
Segment & Unit & TAM & SAM & Platform Rev.\ at SAM \\\midrule
S1 --- B2C marketplace & individuals & \numint{\tamSegOneIndividuals} & \numint{\samSegOneIndividuals} & \$\num[round-precision=1]{\revSamSegOneM}M \\
S2 --- B2B licensing & institutions & \numint{\tamSegTwoInstitutions} & \numint{\samSegTwoInstitutions} & \$\num[round-precision=1]{\revSamSegTwoM}M \\
S3 --- B2B marketplace & companies & \numint{\tamSegThreeCompanies} & \numint{\samSegThreeCompanies} & \$\num[round-precision=1]{\revSamSegThreeM}M \\\midrule
\textbf{Combined} & \textemdash & \textit{not summable} & \textit{not summable} & \textbf{\$\num[round-precision=1]{\revSamAllSegmentsM}M} \\
\bottomrule
\end{tabularx}
\end{table}

\noindent\textit{Units differ by row. The TAM and SAM columns cannot be totalled;
the revenue column can, because dollars are the same unit in every row. ``Platform
Rev.\ at SAM'' is the annual revenue accruing to Great Wall if the entire serviceable
market were captured---an upper bound, not a forecast.}

\subsection{Penetration and What Is Actually in the Projections}

\begin{table}[H]
\centering
\begin{tabularx}{\linewidth}{l Y Y Y l}
\toprule
Segment & Penetration of SAM & Customers at Plan & Annual Platform Rev.\ (exit ARR) & In Y1--Y3 P\&L? \\\midrule
S1 & \num[round-precision=1]{\targetShareSegOne}\% \textit{(derived)} & \numint{\targetSubsSegOne} individuals & \$\num[round-precision=1]{\revPlanSegOneM}M & Yes --- books \$\num[round-precision=1]{\revPlanSegOneActualM}M collected \\
S2 & \numint{\targetShareSegTwo}\% \textit{(illustrative)} & \numint{\targetInstSegTwo} institutions & \$\num[round-precision=1]{\revPlanSegTwoM}M & No \\
S3 & \numint{\targetShareSegThree}\% \textit{(illustrative)} & \numint{\targetCompaniesSegThree} companies & \$\num[round-precision=1]{\revPlanSegThreeM}M & No \\\midrule
\textbf{Combined} & \textemdash & \textit{not summable} & \textbf{\$\num[round-precision=1]{\revPlanAllSegmentsM}M} & \textemdash \\
\bottomrule
\end{tabularx}
\end{table}

\noindent Two disclosures matter more than the numbers themselves:

\begin{itemize}[nosep]
  \item \textbf{S1's penetration is derived, not asserted.} The
    \numint{\targetSubsSegOne} figure is \emph{exactly} the Year~3 subscriber count
    produced by the financial model
    (marketing budget $\div$ CAC, less churn), and \$\num[round-precision=1]{\revPlanSegOneM}M
    is \emph{exactly} \texttt{subARRYearThree}. Market sizing and the P\&L can no
    longer drift apart: the share is computed as
    \numint{\targetSubsSegOne}~/~\numint{\samSegOneIndividuals}, not chosen.
  \item \textbf{S2 and S3 contribute \$0 to the Year~1--3 projections.} The
    \$\num[round-precision=1]{\revUpsideExcludedM}M attributable to them at the
    illustrative penetrations above is upside that is deliberately excluded from
    every revenue, COGS, gross-margin, LTV and valuation figure in this document.
    Their unit economics differ structurally from S1 (see below), and blending them
    into a single ARPU would reintroduce exactly the error this restructure removes.
\end{itemize}

% TODO(SOURCE): the 34.3% figure is arithmetic, but the TENSION it exposes is real
% and unresolved. Consumer-software penetration benchmarks (Coinbase 15%, Stripe 20%,
% LastPass/1Password 10-15%) sit well below it. Either the 15% TAM->SAM assumption is
% too conservative, or the Year 3 marketing ramp is too aggressive. Do not present
% 34.3% as benchmarked - it is not.
\noindent\textbf{Flagged tension.} At \num[round-precision=1]{\targetShareSegOne}\% of
S1 SAM, the Year~3 plan sits \emph{above} the consumer-software penetration benchmarks
cited below. Either the \numint{\samPercentSegOne}\% TAM-to-SAM assumption is too
conservative for this population, or the Year~3 marketing ramp is too aggressive for
the market it addresses. This document does not resolve the tension; it surfaces it.

\subsection{S1 --- Individuals Running Self-Hosted Time-Locks}

\textbf{Unit: individuals. ARPU: \$\num{\arpuSegOneGrossMonthly}/month gross,
\$\num{\arpuSegOneMonthly}/month net to the platform at a
\numint{\markPlatformFeePercent}\% take rate.}

The addressable population is individual, non-institutional holders of more than
\$100{,}000 in Bitcoin who self-custody, and who therefore run a self-hosted time-lock
whose delay is set by their own hardware rather than by the protocol minimum. The
marketplace sells them \emph{speed}: the convenience ratio
$R = C(H_{\max})/C(H_{\text{user}})$ is the factor by which outsourcing collapses a
user's own multi-day solve into the protocol's minimum delay.

\begin{itemize}[nosep]
  \item \textbf{Population: \numint{\tamSegOneIndividuals} individuals}, holding an
    average of \$\numint{\tamSegOneAvgHoldingUsd} for roughly
    \$\numint{\tamSegOneAumBn}bn of protected assets. This replaces the previous
    unsourced \numint{3000000} ``Bitcoin users'' figure and is consistent with the
    companion economics paper.
  \item \textbf{Derivation.} No source publishes a count of \emph{individuals}
    holding more than \$100{,}000 in Bitcoin. Henley~\& Partners count
    \numint{145100} Bitcoin millionaires worldwide\cite{henley2025crypto}; scaling
    down one wealth decade using the address-level decade ratios observable on
    chain---\numint{663495} addresses above \$100{,}000 and \numint{119804} above
    \$1{,}000{,}000\cite{bitinfocharts2026}---yields a band of roughly
    \numint{800000} to \numint{1155000} individuals.
  \item \textbf{Context.} \numint{365000000} people owned Bitcoin at
    end-2025\cite{cryptocom2026owners}, so this population is a fraction of one
    percent of owners---consistent with the extreme wealth concentration reported
    on chain\cite{chainalysis2024,triple2023}. Individuals control the majority of
    circulating supply, with more than \$800bn held in self-custody against roughly
    \$300bn in ETFs and corporate treasuries\cite{river2025custody}.
\end{itemize}

% TODO(SOURCE): 1,100,000 is the TOP of the 800k-1,155k band, not its midpoint. The
% research base case is 950,000. The figure is carried at 1,100,000 only for
% consistency with JUSTIFICATION.tex. Scaling individual-level counts using
% ADDRESS-level decade ratios is an unvalidated methodological step.
% TODO(SOURCE): no self-custody rate by wealth band exists in any source. S1 requires
% self-custody and applies NO haircut for it, so the true core is below 1,100,000.
% TODO(SOURCE): the $300,000 average holding is inherited from JUSTIFICATION.tex and
% could not be independently verified; it implies ~22.8% of circulating supply.
% TODO(SOURCE): the 15% TAM->SAM rate is carried over unchanged from the old flat
% model and nothing was found to validate it for this population.
\noindent\textbf{Unsourced inputs in S1 (flagged, not established):} the choice of
\numint{\tamSegOneIndividuals} rather than the \numint{950000} research base case; the
absence of any self-custody haircut; the \$\numint{\tamSegOneAvgHoldingUsd} average
holding; the \numint{\samPercentSegOne}\% TAM-to-SAM rate; and the
\numint{\subBasicMix}/\numint{\subMediumMix}/\numint{\subProfessionalMix}/\numint{\subGoldenMix}
tier mix that produces the \$\num{\arpuSegOneGrossMonthly} weighted average price.

\subsubsection*{Why the S1 Price Holds}

Three independent tests, in increasing order of strength:

\begin{enumerate}[nosep]
  \item \textbf{Direct comparable.} Casa sells recurring key-management security to
    precisely this buyer at \$250/year (Standard) and \$2{,}100/year
    (Premium)\cite{casa2026pricing}. Great Wall's
    \$\numint{\arpuSegOneGrossMonthly * 12}/year sits inside that proven
    corridor, nearer the entry tier.
  \item \textbf{Assets under protection.} The annual price is
    \numint{\arpuSegOneLoadBps} basis points of the anchor holder's
    \$\numint{\tamSegOneAvgHoldingUsd} position---\emph{below} the 15--25~bps that
    spot Bitcoin ETFs charge\cite{etffees2026}. A holder pays less per year for
    coercion resistance than for the convenience of custody through an ETF, and the
    ETF addresses the wrench attack not at all.
  \item \textbf{Threat value.} At \numint{\arpuSegOneLoadBps}~bps the implied
    break-even is roughly a \num{0.14}\% annual probability of total coercive loss,
    against a risk backdrop in which US border device searches rose to
    \numint{55318} in FY2025, up \num[round-precision=1]{32.4}\% over two years\cite{cbp2025searches},
    with a record \numint{14899} devices searched in the April--June 2025
    quarter\cite{fpf2025border}.
\end{enumerate}

\noindent Pricing is currently derived cost-plus (provider electricity cost, times a
tier markup, grossed up for the platform fee), which has no relationship to the value
of preventing a coerced transfer. For this population cost-plus almost certainly
\emph{under}prices; a migration to value-based pricing denominated in basis points of
protected assets is the natural forward path once real tier-mix data exists.

\subsubsection*{Privacy and Travel: A Wedge, Not a Population}

Border device searches are a low-probability, high-consequence event of exactly the
shape a duress-resistant time-lock addresses: \numint{55318} devices searched in
FY2025, of which \numint{13590} belonged to US citizens\cite{cbp2025searches}, against
roughly \numint{419000000} traveller encounters. No source quantifies the overlap
between international travellers and Bitcoin self-custodians, so this sub-population
is carried at \textbf{zero additional headcount} and used as a conversion and
messaging wedge inside the S1 core.

% TODO(SOURCE): the overlap rate between international travellers and Bitcoin
% self-custodians is COMPLETELY UNSOURCED. Any additive traveller population would be
% a placeholder. It is deliberately carried at +0 here. Do not add one.
% TODO(SOURCE): CBP counts DEVICES and EVENTS, not distinct people. The ~419M
% traveller denominator comes from secondary reporting, not from cbp.gov directly.

\subsubsection*{Journalists and Human Rights Defenders: Mission Reach at Zero ARPU}

Roughly \numint{\segOneMissionReachJournalists} journalists are organised in
professional unions worldwide\cite{ifj2026}; more than \numint{4000} requests for
emergency digital-security assistance reached the Access Now helpline in a single
year\cite{accessnow2025helpline}; \numint{330} journalists were jailed as of
1~December 2025\cite{cpj2026census}. This population has revealed, documented demand
for coercion resistance and \textbf{near-zero ability to pay}. It is modelled at an
ARPU of \$\numint{\arpuSegOneMissionMonthly} as sponsored seats---a marketing or COGS
line, never a revenue line---and is \textbf{not} added to any subscriber count. The
entire global funder pool that could sponsor such seats is small: the Human Rights
Foundation's Bitcoin Development Fund disburses rounds of roughly \$325{,}000 to
\$455{,}000 across 12--26 projects\cite{hrf2025bdf}.

% TODO(SOURCE): the ~600,000 IFJ figure is a UNION MEMBERSHIP count and is a floor,
% not a population; no authoritative count of professional journalists worldwide
% exists. There is no published count of human rights defenders at all.
% TODO(SOURCE): any grant offset for sponsored seats (research suggested $50K-$250K/yr)
% is an inference from observed HRF round sizes. Great Wall has no grant history and
% no source supports that it would receive such funding. No grant revenue is modelled.

\subsection{S2 --- Institutions Licensing Custodial TGPO}

\textbf{Unit: institutions. ARPU: \$\numint{\arpuSegTwoAnnualLicence}/year licence plus
a one-off \$\numint{\arpuSegTwoIntegrationFee} integration fee
(\$\numint{\arpuSegTwoYearOneContract} first-year contract value).
Marketplace revenue: \$\numint{\arpuSegTwoMarketplaceRevenue}.}

In S2 the \emph{institution} imposes the delay on its own customers' outbound
transfers. It follows structurally that S2 generates \textbf{no marketplace revenue
whatsoever}: an institution that wants the delay to be real will never buy solving
capacity to shortcut it. There is no gross transaction volume, no provider side, no
\numint{\markPlatformFeePercent}\% take rate and no payment processing on gross.
S2 economics are licence, seat and integration fees only.

\subsubsection*{The Driver: Brazilian Pix and \textit{Sequestro Rel\^{a}mpago}}

The strongest external validation for the entire product thesis is that the Brazilian
central bank has \emph{already conceded that delay is the countermeasure}, and
mandated it. Resolu\c{c}\~{a}o BCB n.\ 142 caps person-to-person transfers at
R\$1{,}000 in the 20:00--06:00 window and requires a \textbf{minimum 24-hour period}
before an increase to a transfer limit takes effect, while decreases take effect
immediately\cite{bcb2021res142}. That is precisely the asymmetric time-lock semantics
of TGPO---the action that reduces protection is delayed, the action that increases it
is instant---imposed by regulation on every participant. The central bank's president
described the motivating scenario in these terms: a victim of a
\textit{sequestro rel\^{a}mpago} at 02:30 who is forced to make a transfer cannot cancel
it, so transfers in those hours must be deferred to the following
day\cite{cnnbrasil2021pix}.

The ex-post remedy the regulator built instead is measurably failing. Of
R\$\num[round-precision=2]{7.77}bn of \emph{confirmed-fraud} Pix value in 2025, only
\num[round-precision=2]{8.96}\% was actually refunded (2024:
\num[round-precision=2]{8.11}\%), the principal stated reason being insufficient
balance---the money is already gone\cite{bcb2026med}. Recovery after a coerced instant
transfer is roughly \numint{90}\% impossible, which is the entire argument for an
ex-ante time-lock.

The incidence is not marginal. A national victimisation survey found
\num[round-precision=1]{0.8}\% of Brazilians aged 16+ reported suffering a
\textit{sequestro rel\^{a}mpago} in the prior twelve months---\numint{1267167} people, up
from \num[round-precision=1]{0.6}\% the previous edition---and identified coercion as
the mechanism linking physical and digital crime, since the offender must now compel
the victim to surrender credentials rather than simply take a
device\cite{fbsp2025vitimizacao}. \numint{830890} mobile phones were stolen or robbed
in Brazil in 2025, of which \numint{308723} were \textit{roubos}---taken by violence or
threat, with the victim present and coercible\cite{fbsp2026anuario}.

% TODO(SOURCE): the 1,267,167 figure is a SURVEY-BASED POPULATION PROJECTION from
% self-reports, not a count of police reports. Official recorded counts are far lower
% and were not retrieved. Expect an investor to push on that gap.

\subsubsection*{Institution Counts}

\begin{itemize}[nosep]
  \item \textbf{Brazil: \numint{\tamSegTwoInstBrazil} active Pix participants} at
    December 2025, up \num[round-precision=1]{5.6}\% year on year---a hard,
    annually republished, regulator-published count of exactly the institutions that
    could embed custodial TGPO\cite{bcb2026pixreport}.
  \item \textbf{Colombia: \numint{\tamSegTwoInstColombia}} entities linked to Bre-B,
    the interoperated instant-payment system\cite{banrep2026breb}.
  \item \textbf{Argentina: \numint{\tamSegTwoInstArgentina}} registered payment
    service providers offering payment accounts\cite{bcra2025pagos}.
  \item \textbf{Mexico: \numint{\tamSegTwoInstMexico}}, being 50 authorised banks plus
    89 authorised Financial Technology Institutions\cite{cnbv2026banca}.
\end{itemize}

\noindent SAM is cut to \numint{\samSegTwoInstitutions} independent buyers
(\num[round-precision=1]{\samPercentSegTwo}\% of TAM) because the Pix participant
count includes several hundred individual credit cooperatives that run on shared
core-banking platforms and do not procure independently. Brazil alone contributes
\numint{\samSegTwoInstBrazil} and is the beachhead.

% TODO(SOURCE): the claim that credit cooperatives buy through shared centrals
% (Sicoob, Sicredi, Ailos, Cresol, Unicred) rather than individually is UNSOURCED, and
% it is what removes ~500 institutions from the 926. It materially drives the SAM.
% TODO(SOURCE): the Brazil figure of 400 is arithmetic over three separately sourced
% components (~141 banks in BCB's STR list, >200 authorised payment institutions,
% ~60 SCD/SCFI) whose OVERLAP could not be verified. Only the 926 is a hard count.
% TODO(SOURCE): BCB does not publish head counts of the 926 by segment. Chase the
% IF.data / "Relacao de Instituicoes em Funcionamento" registry before this goes to an
% investor. Peru, Chile, Ecuador and Central America were not sized at all.

\subsubsection*{Licence Pricing}

The \$\numint{\arpuSegTwoAnnualLicence} blend is anchored deliberately \emph{below}
the two hardest available benchmarks. BioCatch---the closest true comparable, selling
a single behavioural anti-fraud control to banks---finished 2025 with more than
\$185M ARR across \numint{340} financial-institution customers, a mean of roughly
\$\numint{544000} per institution\cite{biocatch2026}. At the top end, Caixa
Econ\^{o}mica Federal's disclosed emergency contract for cloud facial-biometric
authentication was R\$\num[round-precision=1]{8.1}M for 180
days\cite{caixa2026biometria}, annualising to several million dollars for one control
at one bank. Great Wall prices at roughly one-fifth of the BioCatch mean because
custodial TGPO is one narrow, verifiable control---an enforced,
cryptographically non-shortcuttable delay on the specific class of outbound transfer a
coerced customer can be forced to authorise---not a fraud platform with models, case
management and analyst tooling. It should price as a module, not a platform
replacement.

Two arguments support a defensible floor. First, the compliance hook: the regulator
already mandates the mechanism\cite{bcb2021res142}, so the institution is buying a
better implementation of an existing legal obligation, which prices more reliably
than discretionary security spend. Second, loss avoidance: with roughly
\num[round-precision=0]{9}\% of confirmed-fraud value recovered\cite{bcb2026med},
system-wide unrecovered confirmed-fraud value ran in the billions of reais in 2025,
against which a \$\numint{\arpuSegTwoAnnualLicence} licence is trivial for any
institution carrying even a small share of that exposure.

% TODO(SOURCE): NO public benchmark exists for what a bank pays for a coerced-transfer
% or enforced-delay control specifically. Every benchmark found is for a broader
% product. The step from BioCatch's $544K mean to $120K is a scope judgement.
% TODO(SOURCE): the $75,000 integration fee has no sourced comparable at all.
% TODO(SOURCE): the 10/100/290 tier split and $900K/$220K/$60K price points that blend
% to $121K are constructed, not observed.
% TODO(SOURCE): no evidence that any Brazilian institution today carries a budget line
% for coerced-transfer delay as a discrete product. Willingness to pay is assumed.
% TODO(SOURCE): the 8% five-year penetration is illustrative; no adoption curve for a
% regulated-payments security module in LatAm was sourced.
\noindent\textbf{Unsourced inputs in S2 (flagged, not established):} the scope
discount from the BioCatch mean to \$\numint{\arpuSegTwoAnnualLicence}; the
\$\numint{\arpuSegTwoIntegrationFee} integration fee; the internal tier ladder that
blends to \$\numint{\arpuSegTwoBlendCheck}; the cooperative-procurement haircut; and
the \numint{\targetShareSegTwo}\% penetration rate.

\subsection{S3 --- Companies Buying Solving Capacity for Digital Assets}

\textbf{Unit: companies. ARPU: \$\numint{\arpuSegThreeGrossMonthly}/month gross,
\$\numint{\arpuSegThreeMonthly}/month net to the platform.}

S3 sells outsourced puzzle-solving capacity to businesses protecting \emph{digital}
assets: self-custodied treasury keys, signing credentials, vault and root secrets.
Cryptographic delay is necessary only where all three of the following hold, and it is
this filter---not the size of the retail sector---that determines the segment:

\begin{enumerate}[nosep]
  \item the protected asset is a \textbf{secret}, not a physical object, so no
    enclosure can gate access to it;
  \item the adversary may be an \textbf{insider} holding legitimate override
    authority, so an override-capable lock is worthless;
  \item the adversary may be the \textbf{vendor or custodian}, so no trusted third
    party may hold a bypass.
\end{enumerate}

\noindent The documented failure modes are exactly these. Private-key compromise
accounted for \num[round-precision=1]{43.8}\% of all crypto value stolen in 2024, the
single largest attack vector, out of \$\num[round-precision=1]{2.2}bn
stolen\cite{chainalysis2025hacks}. The largest exchange theft on record worked by
compromising a multisig \emph{signing interface} rather than the keys, and a second
nine-figure loss drained a multisig operated under a third-party custody
arrangement\cite{wazirx2024hack}. No safe, timer or custodian override addresses any
of those; a cryptographic delay does.

\begin{itemize}[nosep]
  \item \textbf{TAM: \numint{\tamSegThreeCompaniesLow}--\numint{\tamSegThreeCompaniesHigh}
    companies worldwide} (modelled at \numint{\tamSegThreeCompanies}), dominated by
    licensed crypto-asset service providers---more than \numint{7360} tracked across
    75 countries\cite{thebanks2025vasp}.
  \item \textbf{SAM: \numint{\samSegThreeCompanies} companies}, the Brazil-anchored
    beachhead. \numint{12053} Brazilian organisations declared crypto on their balance
    sheets\cite{receita2022cripto}, of which only the self-custodying subset is
    addressable---just \num[round-precision=1]{7.6}\% of businesses fully self-custody,
    and fewer than \numint{100} corporate bitcoin treasuries of meaningful size exist
    anywhere\cite{river2025business}.
  \item \textbf{Pricing} is bracketed between the S1 Golden tier
    (\$\numint{\subGoldenPrice}/month, floor: a company runs more puzzles, more
    signers and longer delays than one individual) and the verified
    \$\num[round-precision=2]{435.56}/month Loomis SafePoint contract price
    (ceiling---and that bundles armoured transport, cash counting and a loss
    guarantee)\cite{loomis2016safepoint}. Model as a range of
    \$\numint{\arpuSegThreeGrossMonthlyLow}--\$\numint{\arpuSegThreeGrossMonthlyHigh}
    gross per month, not a point estimate.
\end{itemize}

\noindent\textbf{S3 does not scale with site count.} Treasury keys are held centrally,
so a jewellery chain with forty stores buys one time-lock configuration, not forty---%
unlike a time-delay safe, which is bought per site. Modelling S3 per-site would
overstate its revenue by an order of magnitude. S3 is therefore a \textbf{high-ARPU,
low-count, strategic} segment that validates the technology with demanding buyers and
generates reference customers. It is not the volume engine; S1 is.

% TODO(SOURCE): the product category HAS NO MARKET. No business has ever paid for
% outsourced time-lock-puzzle solving, so every S3 ARPU figure is constructed by
% analogy to adjacent products. There is no comparable transaction anywhere.
% TODO(SOURCE): the 28% take rate is carried over from S1 with no B2B benchmark.
% TODO(SOURCE): applies River's US 7.6% full-self-custody rate to Brazilian companies;
% no Brazil-specific corporate custody split exists.
% TODO(SOURCE): the 6,000 upper bound reads River's unquantified "most use a
% third-party + self-custody hybrid" as roughly 50%. River publishes no such number.
% TODO(SOURCE): the Receita Federal count (12,053) is from August 2022 and has NOT
% been extrapolated forward, despite Brazil's VASP regime taking effect in 2026.
% TODO(SOURCE): no conversion or penetration rate from addressable population to
% paying customers can be estimated for a category that does not exist; the 10% used
% is illustrative. No LatAm business counts outside Brazil were sourced.
\noindent\textbf{Unsourced inputs in S3 (flagged, not established):} the
\$\numint{\arpuSegThreeGrossMonthly} price point; the
\numint{\markPlatformFeePercent}\% take rate applied to a B2B buyer; the transfer of a
US self-custody rate to Brazil; the upper bound of the SAM range; the currency of the
2022 Brazilian corporate-crypto count; and the
\numint{\targetShareSegThree}\% penetration rate.

\subsection{Existing-Market Comparable: Time-Delay Safes}

The nearest thing to an existing market for this product is the \textbf{time-delay
safe}: a mature, regulated, commercially successful product sold to convenience
stores, pharmacies, fuel stations, supermarkets and lottery retailers, whose entire
value proposition is that access is deliberately slowed and that the slowdown is
\emph{advertised on the door}. It is the proof that a delay is a saleable good.

\begin{itemize}[nosep]
  \item \textbf{The deterrent effect is documented by regulators, not vendors.} After
    Alberta mandated time-delayed safes in pharmacies, robberies fell from
    \numint{89} to \numint{14} in Calgary and from \numint{39} to zero in Edmonton;
    British Columbia recorded a \numint{94}\% decrease\cite{ocp2024safes}. Ontario
    reported an \numint{82}\% decrease after full implementation across all community
    pharmacies\cite{oacp2024safes}. A national US pharmacy chain reported a
    \numint{50}\% decrease across 45 states\cite{cvssafes2015}.
  \item \textbf{The effect depends on publishing the mechanism.} Every source
    attributes the result to prominent signage plus sector-wide coverage, so that
    offenders cannot shop for non-adopters\cite{ocp2024safes,cvssafes2015}. Florida
    law goes further and \emph{requires} a conspicuous notice at the entrance
    advertising that cash is limited\cite{flsenate812173}. This is direct empirical
    support for Great Wall's anti-obscurity posture: the mechanism is published
    precisely because publication is what deters.
  \item \textbf{Businesses pay recurring fees for delay.} A publicly filed Loomis
    SafePoint agreement prices delay-plus-service at
    \$\num[round-precision=2]{435.56} per month per safe unit on a five-year
    term\cite{loomis2016safepoint}.
  \item \textbf{The category already exists in Brazil.} \textit{Cofres com retardo}
    with configurable 1--99 minute delays are sold to fuel stations, pharmacies and
    other commercial establishments explicitly to reduce losses in
    \textit{assaltos rel\^{a}mpagos}\cite{importec2024retardo}; smart safes are adopted
    by supermarkets, pharmacies, lottery retailers and fuel
    stations\cite{prosegur2024cofres}. The broader safes-and-vaults market is
    estimated in the single-digit billions of
    dollars\cite{grandview2024,mordor2024}, and private protective services are a
    separate multi-billion-dollar market serving the same
    fear\cite{alliedmarket2023}.
  \item \textbf{But the buyer population is smaller than it looks.} Brazilian
    official statistics count \numint{10607110} companies and organisations, of which
    only \numint{80142} have 50--249 employees---and in \emph{all} of retail only
    \numint{7251} enterprises are that size\cite{ibge2024cempre}. Any model assuming
    tens of thousands of mid-size Brazilian retail buyers is wrong.
\end{itemize}

% TODO(SOURCE): NO published market size exists for the TIME-DELAY segment
% specifically, anywhere. The safes-and-vaults totals are proxies from secondary
% research vendors whose 2025 estimates disagree by roughly 2x ($4.8bn to $9.13bn).
% TODO(SOURCE): no source was found quantifying an insurance premium discount for
% installing a time-delay safe. Vendors assert underwriting benefits without evidence.
% Do not put a number on it.
% TODO(SOURCE): the Loomis price is a single 2016 US municipal contract with a CPI
% escalator, not a market average; the current equivalent is higher by an unknown amount.
\noindent\textbf{Where the analogy stops, stated plainly.} Every deterrence statistic
above was produced by cheap mechanical and electronic timers, not by any cryptographic
construction. A business protecting cash or jewellery in a steel box has no use for
sequential squaring: it buys a timer lock for a few hundred dollars, bolts it on, and
is done. Those numbers justify the \emph{concept} of advertised delay; they do not by
themselves establish a market for \emph{cryptographic} delay. That market exists only
where the protected asset is a secret rather than an object---which is why S3 is scoped
to digital assets and is deliberately small.

\subsection{Market Share Benchmarks}
The S1 target share is derived from the financial model rather than asserted (see
above). The following comparables are offered as context for \textbf{S1 only}:

\begin{itemize}[nosep]
  \item \textbf{Consumer penetration comparables:} Coinbase's 15\% crypto exchange
    capture\cite{coinbase2021}, Stripe's 20\% payment processing
    share\cite{stripe2023}, and LastPass/1Password's 10--15\% password management
    penetration\cite{lastpass2023}.
  \item \textbf{These do not transfer to S2 or S3.} Consumer-software penetration
    curves say nothing about enterprise licence adoption in regulated payments (S2)
    or about a B2B category that does not yet exist (S3). No benchmark is claimed for
    either; their penetration rates are labelled illustrative in the table above.
  \item \textbf{Strategic advantages:} Partnership with a market leader provides
    distribution, brand credibility and accelerated acquisition, typically doubling
    organic growth rates\cite{reforge2022,hagiu2015}.
\end{itemize}

\subsection{Market Dynamics}
\begin{itemize}[nosep]
  \item \textbf{Bitcoin Adoption:} Growing mainstream adoption drives demand for
    security tools
  \item \textbf{Self-Sovereignty Trend:} ``Not your keys, not your coins'' philosophy
    expanding the S1 base
  \item \textbf{Regulatory Tailwind:} Central banks are already mandating delay as an
    anti-coercion control\cite{bcb2021res142}, which converts S2 from discretionary
    security spend into a compliance line
  \item \textbf{Privacy Concerns:} Increasing demand for anonymous computation
    services
  \item \textbf{Underserved Market:} Limited competition in the anonymous marketplace
    segment
\end{itemize}

\subsection{2025 Market Inflection}
\label{sec:market-inflection}
The year 2025 witnessed three seemingly unrelated developments:
\begin{enumerate}
  \item A sharp rise in reported wrench attacks targeting cryptocurrency holders\cite{lopp2024,trmlabs2025,theblock2025,chainalysis2025}---a trend that continues to the time of writing;
  \item The breaking of Bitcoin's historical 3-up-1-down price cycle, with Bitcoin closing the year at a loss in what should have been a positive year\cite{cnbc2025btc,ccn2025btc};
  \item A sustained bull market in precious metals\cite{bloomberg2025gold,jpmorgan2025gold}.
\end{enumerate}

While the causal links between these events remain unproven, a plausible narrative emerges: both large and small crypto holders, increasingly viewing ownership as a security liability, may be diversifying into physical assets. The reasoning is straightforward: if holding cryptocurrency incurs substantial security costs (vaults, alarms, surveillance, armed protection) and usage friction (displacement to sites of custody, extra authentication steps, reduced portability), one might as well bear those costs for precious metals instead, thereby eliminating the technological complexity inherent to crypto custody.

Notably, any such capital flight would likely remain invisible in public data. Holders motivated by security concerns would execute these transitions discretely precisely because, to them, security through obscurity appears paramount given the gravity and persistence of the wrench attack threat. This suggests the observable trend may significantly understate the actual magnitude of the shift. Additionally, that explanation also suggests an \href{https://en.wikipedia.org/wiki/Anosognosia}{\textit{anosognostic}} mechanism: the problem, by its nature, causes individuals to prefer, for their own security, not to talk about it, which impedes accurate diagnose and potential treatments --- the exact conditions for ``you didn't know you needed it'' innovations.

This dynamic represents a significant tailwind for Great Wall. Rather than abandoning cryptocurrency, holders can retain the benefits of digital assets while neutralizing the security concerns that would otherwise drive them toward traditional stores of value.

\subsection{Competitive Landscape}
\begin{itemize}
  \item \textbf{Direct Competition:} Limited due to anonymous marketplace complexity
  \item \textbf{Indirect Competition:} Traditional cloud computing lacks privacy features; time-delay safes solve the physical case only and cannot gate a secret
  \item \textbf{Barriers to Entry:} Trust and reputation system creates moat
  \item \textbf{First-Mover Advantage:} Early provider network difficult to replicate
\end{itemize}
